Agentic software systems are increasingly expected to invoke tools, consume budgets, and produce externally consequential outputs, yet most current runtimes still expose execution primarily through logs and traces. Those artifacts are useful for observability, but they do not by themselves provide a bounded, reviewable, and independently checkable account of what was intended, what was authorized, what was executed, and what evidence was exported. This paper frames that mismatch as an execution evidence problem rather than a logging problem.

We present an execution evidence architecture for agentic software systems and instantiate it as a reproducible reference path in an artifact package omitted here for double-blind review. The architecture links five stages of runtime semantics: persona projection, intent/action/result objects, governance checkpointing, execution-integrity validation, and bounded audit receipts. The controlled evaluation directly tests the latter four evidence-preserving stages through a deterministic 16-task suite, a unified five-file export contract, a third-party bundle review script, minimal live CrewAI and LangChain baselines, paired same-framework comparisons in both frameworks with and without the evidence layer, four architectural ablations, and a small falsification suite for negative controls.

Across the checked-in evaluation artifacts, the full evidence-chain configuration improves average explicitness from 1.0 to 5.0, replayability from 3.0 to 5.0, tamper sensitivity from 0.0 to 4.81, audit boundedness from 2.0 to 5.0, and evidence-stage exposure from 1.0 to 5.0 relative to the local trace-oriented baseline. The minimal live CrewAI and LangChain baselines each reach replayability 4.0, but leave the other evidence dimensions at observability-like levels. In both same-framework paired comparisons, wrapping the identical live runtime with the evidence layer raises the corresponding evidence metrics to 5.0, 5.0, 4.81, 5.0, and 5.0. Ablation results show targeted degradation when intent capture, governance persistence, integrity validation, or receipt export is removed. Under an eight-scenario falsification suite, the updated review contract rejects cross-run replay mismatches, cross-bundle receipt swaps, receipt-digest forgeries, payload tampering, false tamper claims, and policy-bypass claims, yielding a 1.0 detection rate across all derived bundles and showing that the contract is not limited to field-presence checks. An appendix-only blinded review kit is included as workflow scaffolding, but not used as empirical evidence. In this controlled harness, the results support a bounded claim that observability surfaces and verifiability surfaces are not interchangeable.
